%==============================================================================
% Poster graduaatsproef
%==============================================================================

\documentclass[a0,portrait]{hogent-poster}

% Sampled from graphics/forcebit_logo.jpeg so the poster headings match the project branding.
\definecolor{forcebit-blue}{RGB}{82,107,138}
\colorlet{title}{forcebit-blue}

% Info over de opleiding
\course{Graduaatsproef}
\studyprogramme{Graduaat in het Programmeren}
\academicyear{2025-2026}
\institution{Hogeschool Gent, Valentin Vaerwyckweg 1, 9000 Gent}

% Info over de graduaatsproef
\title{FORCEBIT TICKETING}
\subtitle{Een support ticketing applicatie met een gelaagde\\architectuur}
\author{Alex Delvoye}
\email{alexdelvoye@outlook.be}
\supervisor{Emlin Oguz Inci}
\cosupervisor{Sacha De Maere (Forcebit)}

\specialisation{Programmeren}
\keywords{Ticketing, ASP.NET Core, React Native, MySQL, onderhoudbaarheid}
\projectrepo{https://github.com/alexdelvoye/GraduaatsProef_TicketingSystem.git}

\begin{document}

\maketitle

\begin{abstract}
Forcebit Ticketing is een proof of concept voor een centrale supportworkflow.
De applicatie brengt klantvragen, antwoorden, statussen en bijlagen samen in tickets.
De focus ligt op een onderhoudbare opbouw: een gelaagde backend, duidelijke
domeinregels en een frontend met herbruikbare schermlogica.
\end{abstract}

\begin{multicols}{2}
\small

\section{Probleem}

Bij Forcebit verliep klantcommunicatie vooral via telefoon en e-mail. Daardoor
raakten antwoorden, screenshots, bijlagen en statussen verspreid over verschillende
kanalen. De vraag werd daarom hoe ik een onderhoudbare ticketingapplicatie kon
bouwen waarin klant en administrator dezelfde duidelijke supportworkflow volgen.

\begin{center}
  
\includegraphics[width=0.25\linewidth]{forcebit_logo}
\end{center}

\section{Oplossing}

Elke supportvraag wordt centraal bewaard als ticket. De eerste beschrijving van de
klant wordt meteen het eerste bericht in de conversatie. Daarna kunnen klant en
administrator antwoorden en bijlagen toevoegen zolang het ticket open blijft.

\textbf{Workflow}
\begin{itemize}
  \item \texttt{New}: support heeft nog niet geantwoord.
  \item \texttt{Open}: de conversatie loopt of het ticket werd heropend.
  \item \texttt{Closed}: het ticket is afgehandeld.
\end{itemize}

Een heropend ticket wordt opnieuw \texttt{Open}, niet \texttt{New}, omdat het al een
geschiedenis heeft.

\section{Gelaagde backend}

De backend is gebouwd met ASP.NET Core en opgesplitst in vier lagen:

\begin{itemize}
  \item \textbf{API}: HTTP-verzoeken, JWT-claims en requestdata.
  \item \textbf{Services}: use cases zoals antwoorden en statussen wijzigen.
  \item \textbf{Persistence}: EF Core, repositories, configuraties en migraties.
  \item \textbf{Domain}: entiteiten, enums en domeinregels.
\end{itemize}

Deze verdeling houdt controllers dun. Toegang, statusovergangen en antwoorden op
gesloten tickets worden niet in schermen of controllers beslist, maar centraal
bewaakt door services en domeinregels.

\begin{center}
  \captionsetup{type=figure}
  \includegraphics[width=0.82\linewidth]{Backend_architecture}
  \captionof{figure}{Backendlagen en belangrijkste afhankelijkheden.}
\end{center}

\section{Frontend en gebruiksgemak}

De frontend is gebouwd met Expo React Native. Schermen tonen de pagina's, hooks
beheren gedrag, API-modules communiceren met de backend en componenten houden de
UI herbruikbaar. Dezelfde paginatiehook wordt gebruikt voor gesprekken, klanten en
ticketgroepen.

\textbf{Belangrijke UI-keuzes}
\begin{itemize}
  \item zoeken, filters, statusgroepen en paginatie;
  \item chatbubbels met onderscheid tussen klant en admin;
  \item preview voor PNG/JPG-afbeeldingen;
  \item uploadfeedback voor de 20 MB selectielimiet;
  \item e-mailnotificaties via Brevo bij adminacties.
\end{itemize}

\begin{center}
  \captionsetup{type=figure}
  \includegraphics[width=0.82\linewidth]{Ticket_conversation}
  \captionof{figure}{Ticketdetail met chatbubbels en afbeeldingspreview.}
\end{center}

\section{Onderhoudbaarheid}

De belangrijkste technische keuze is Single Responsibility. Controllers behandelen
HTTP, services voeren workflows uit, repositories halen data op en domeinregels
bewaken businessregels. In de frontend geldt dezelfde gedachte: schermen, hooks,
API-modules, validatie en componenten blijven gescheiden.

\section{Resultaat en vervolg}

Het resultaat is een werkend proof of concept voor een realistische supportworkflow:
registreren, inloggen, tickets aanmaken, antwoorden, bijlagen toevoegen, zoeken,
filteren, pagineren, sluiten, heropenen en notificeren.

\textbf{Logische vervolgstappen}
\begin{itemize}
  \item productiegerichte opslag voor bijlagen, zoals object storage;
  \item administrators koppelen aan categorieën, onderwerpen of tickets;
  \item server-side paginatie en automatische tests;
  \item verdere uitwerking richting mobiele app.
\end{itemize}

\end{multicols}
\end{document}
